
\subsubsection{6.1 Key Findings}

The experiments reveal three critical findings:

\textbf{Finding 1: RLHF datasets contain verifiable safety violations.} The HH-RLHF harmless subset achieves 45.6\% genuine violation base-rate (95\% CI: [41.3\%, 50.0\%]), validating dataset quality. This addresses a gap in RLHF literature—prior work assumes preference labels are valid but does not empirically audit base-rates. The validation protocol provides a framework for dataset quality assessment.

\textbf{Finding 2: Human annotators achieve substantial consistency with explicit criteria.} Inter-annotator agreement (κ=0.724) demonstrates that violation detection is learnable and reproducible. This validates that human-detectable structure exists in preference data—annotators apply explicit safety criteria with 83.6\% alignment to original labels.

\textbf{Finding 3: Pretrained embeddings fail to capture alignment structure despite human detection consistency.} RoBERTa embeddings show effectively random separation (Cohen's d=0.034, only 8.5× above baseline), revealing a disconnect: human-detectable structure (κ=0.724) does not equal embedding-space structure (d=0.034). With 160,000 samples providing statistical power exceeding 0.99, this is not a data insufficiency issue but a representation limitation.

\textbf{Synthesis:} Pretrained encoders optimize for general semantic similarity through masked language modeling rather than safety-specific discrimination. Alignment violations span diverse semantic content that occupies overlapping embedding regions despite human-detectable differences. Safety distinctions operate on semantic dimensions orthogonal to general-purpose pretraining objectives.

\subsubsection{6.2 Limitations}

Several limitations bound the scope of these findings:

\textbf{L1: Single-encoder negative result.} Only RoBERTa-base was tested. Safety-fine-tuned encoders or reward model embeddings may capture structure that general-purpose encoders miss. RoBERTa is a widely-used baseline that establishes standard pretrained models are insufficient. Multi-encoder validation and safety-specialized alternatives remain as future work. If any encoder shows d≥0.5, it would refine the conclusion from "embeddings fail" to "general-purpose embeddings fail."

\textbf{L2: Annotation consistency used untrained data.} H-M1 achieved κ=0.724 using h-e1 annotators without the full training protocol, potentially underestimating ceiling performance. This is a proof-of-concept demonstrating analysis infrastructure. The κ=0.724 baseline already shows substantial agreement (exceeds 0.70 threshold). Prior annotation studies suggest training improves κ by 0.10-0.15, which would strengthen rather than weaken the conclusions. Future work should recruit trained annotators.

\textbf{L3: HH-RLHF harmless subset only.} Results are specific to safety violations in conversational AI. Helpfulness preferences, other RLHF datasets, and non-English data remain untested. Focused scope enables controlled experiments. The harmless subset isolates safety-related judgments, directly testing the hypothesis. The validation protocol is dataset-agnostic and generalizable.

\textbf{L4: Geometric structure hypotheses untested.} H-M3 (severity-distance correlation) and H-M4 (encoder invariance) were blocked by H-M2 failure. Testing these requires clustering to exist first. Without baseline separation (d=0.034), analyzing cluster properties is not meaningful. If alternative encoders show clustering (d≥0.5), these tests become viable.

\subsubsection{6.3 Broader Impact}

This work advances AI safety through three mechanisms:

\begin{enumerate}
\item \textbf{Dataset quality validation:} The base-rate protocol (45.6\% genuine violations) provides empirical grounding for RLHF dataset quality, enabling researchers to verify annotation authenticity.
\end{enumerate}

\begin{enumerate}
\item \textbf{Methodological contribution:} The three-hypothesis cascade (existence → consistency → structure) offers a framework for testing complex causal chains in alignment research.
\end{enumerate}

\begin{enumerate}
\item \textbf{Research redirection:} By demonstrating pretrained encoder insufficiency, this work redirects effort toward safety-specialized representations and reward model mechanistic interpretability.
\end{enumerate}

No significant concerns are identified. This work is defensive (alignment evaluation methodology) rather than offensive. The negative finding limits potential misuse—demonstrating that general-purpose embeddings do not capture safety structure reduces risk of adversarial exploitation via embedding-space attacks.

Scientific progress requires both positive and negative results. By testing and refuting the geometric manifold hypothesis for pretrained embeddings, this work establishes knowledge boundaries: what does not work, why it fails, and what alternatives to pursue.

\subsubsection{6.4 Theoretical Interpretation}

Why do pretrained embeddings fail despite human detection consistency? Three complementary explanations:

\textbf{Explanation 1: Optimization objective mismatch.} RoBERTa optimizes masked language modeling (predict masked tokens from context). Safety violations are not defined by missing tokens but by content violating policies. A response can be grammatically coherent, semantically plausible, and still unsafe. Pretrained embeddings capture linguistic plausibility rather than normative safety.

\textbf{Explanation 2: Semantic diversity of violations.} Toxicity, harmful instructions, and misinformation occupy distant semantic regions despite unified "unsafe" classification. Unlike clustering tasks where semantic similarity approximates task similarity (e.g., sentiment analysis), safety violations lack semantic cohesion. They are defined by violating policies rather than sharing linguistic features.

\textbf{Explanation 3: Implicit versus explicit features.} Human annotators detect violations using explicit criteria (policy violations, harm potential). Pretrained embeddings encode distributional statistics (word co-occurrence, syntactic patterns). These feature spaces are non-overlapping for safety tasks—the signals humans use are not manifest in pretraining objectives.

\textbf{Supporting evidence:} The near-zero effect size (d=0.034) with massive sample size (160,000) rules out weak-but-real effects. This is a fundamental representation failure rather than a statistical power issue. If safety structure existed in pretrained embeddings, 160,000 samples would reveal it.
